\documentclass[12pt]{article}
\usepackage[a3paper, margin=1in]{geometry}
\usepackage{amsmath, amssymb, amsthm, ulem, booktabs}

\begin{document}

\section*{Domácí úkol 2.2}

Vypracovali: Daniel \textit{``Localhost"} Ransdorf, Jan \textit{``kokeshh"} Kodeš \hfill Podpis: \rule{4cm}{0.4pt}


\bigskip

Označme prvky matice $A$:

\[
  \left(\begin{array}{cc}
    a & b \\ b & d
  \end{array}\right) := A
\]

Zvolme dva vektory závislé na proměnné (úhlu) $x\in\mathbb{R}$, které vzniknou rotací kanonické báze.


\[
  \left(\vec{v}_1, \vec{v}_2\right) := \left(
    \left(\begin{array}{c}
      \cos(x) \\ \sin(x)
    \end{array}\right),
    \left(\begin{array}{c}
      -\sin(x) \\ \cos(x)
    \end{array}\right)
  \right)
\]

Tyto vektory tvoří ortogonální posloupnost vzhledem k std. sk. s. nezávisle na $x$, ověřme:

\[
  \vec{v}_1 \cdot \vec{v}_2 = -\sin(x)\cos(x) + \sin(x)\cos(x) = 0 \quad \checkmark
\]

Dokažme, že pro libovolný součin ze zadání $<,>$, existuje $x$ takové, že $(\vec{v}_1,\vec{v}_2)$ je
ortogonální i vzhledem k $<,>$.

Jest zvolme libovolnou pozitivně definitní matici $A \in \mathbb{R}^{2x2}$,
rozepišme rovnici kolmosti $\vec{v1},\vec{v2}$:

\begin{align*}
  0 = \left<\vec{v}_1, \vec{v}_2\right>
  &= \left(\begin{array}{cc}
      \cos(x) & \sin(x)
    \end{array}\right)
    \left(\begin{array}{cc}
      a & b \\ b & d
    \end{array}\right)
    \left(\begin{array}{c}
      -\sin(x) \\ \cos(x)
    \end{array}\right) \\
  &= -a\cdot \sin(x)\cos(x) + b\cdot \cos^2(x) - b\cdot \sin^2(x) + d\cdot \sin(x)\cos(x) \\
  &= (d-a)\cdot \sin(x)\cos(x) + b\cdot(\cos^2(x) - \sin^2(x)) \\
  &= (d-a)\cdot \frac{\sin(2x)}{2} + b\cdot \cos(2x)
\end{align*}

Matice byla volena libovolně. Čili ukažme, že existuje úhel $x$, který splňuje požadovanou rovnost. Rozdělme dva případy:

\begin{itemize}
  \item $a-d \neq 0$:
  Získáváme:
  
  \begin{align*}
    0 &= (d-a)\cdot \frac{\sin(2x)}{2} + b\cdot \cos(2x) \\
    (a-d)\cdot \frac{\sin(2x)}{2} &= b\cdot \cos(2x) \\
    \frac{\sin(2x)}{\cos(2x)} &= \frac{2b}{a-d} \quad\quad\quad ,x\neq \frac{\pi}{4}(1+2k) \quad (*)\\
    tan(2x) &= \frac{2b}{a-d} \\
    2x &= arctan\left(\frac{2b}{a-d}\right) \\
    x &= \frac{1}{2} arctan\left(\frac{2b}{a-d}\right) \quad \\
    \text{Ověřme, že řešení }&\text{není zakázaná hodnota }x \text{ dle }(*)\\
    H(arctan) = \left(\frac{-\pi}{2}, \frac{\pi}{2}\right) \quad \Rightarrow \quad x &\in \left(\frac{-\pi}{4}, \frac{\pi}{4}\right)
    \quad \Rightarrow \quad x \neq \frac{\pi}{4}(1+2k) \quad \checkmark
  \end{align*}

  \item $a-d = 0$:
  
    \begin{align*}
      0 &= (d-a)\cdot \frac{\sin(2x)}{2} + b\cdot \cos(2x) \\
      0 &= b\cdot \cos(2x) \\
      \text{Pozorování: pro }b=0\text{ je řešením jakékoli }&x\text{, předpokládejme BÚNO, že }b\neq0 \\
      \Rightarrow \quad 0 &= \cos(2x) \\
      2x &= \frac{\pi + 2k\pi}{2} \quad ,k\in\mathbb{Z} \\
      x &= \frac{\pi + 2k\pi}{4} \quad ,k\in\mathbb{Z} \\
      x_1 &= \frac{\pi}{4}
    \end{align*}

\end{itemize}

Tedy pro libovolně volenou pozitivně definitní matici $A$ jsme našli úhel $x$, se kterým vychází skalární součin $<\vec{v_1},\vec{v_2}> = 0$.
Posloupnost $(\vec{v_1},\vec{v_2})$ tvoří bázi $\mathbb{R}^2$ kolmou vzhledem k std. sk. s. a ted' víme, že i vzhledem k $<,>$.
Tím je tvrzení dokázáno.

\end{document}
